\histitem{算子代数}

我们已经看到，里斯在1913年阐释希尔伯特学派工作时，赋予希尔伯特空间自同态代数 \(\mathcal{L}(\mathrm{E})\) 的作用，也赋予全纯函数演算和谱的抽象概念以作用。在1916年的回忆录中，他似乎已经预示赋范代数是谱理论的自然框架；经过巴拿赫及其学派在20世纪20年代

的工作，这一思想占据中心地位也就不足为奇了。M. 长友 {[}48{]} 于1936年引入巴拿赫代数概念，而 S. 马祖尔 {[}44{]} 于1938年证明，唯一作为域的 C 上赋范代数就是 C 本身（I，第26页，推论2）。

大约在20世纪30年代中期，与谱理论相关，抽象考察希尔伯特空间上的算子代数的机会不断增多。F. 默里和 J. 冯·诺伊曼在一系列深刻且有影响的工作 {[}47{]} 中，系统研究了等于自身双交换子的 \(\mathcal{L}(\mathrm{E})\) 的对合子代数（“冯·诺伊曼代数”）。另一方面，S. 斯廷于1936年提出引入“算子代数”的公理化概念，并深入研究相关结构 {[}70{]}。此外，斯通 {[}74{]} 受谱投影研究启发，于1937年研究其所有元素均为幂等元的含单位代数（“布尔代数”：参见 ÉHM，第146页）。他注意到，若 X 是局部紧拓扑空间，则开集的特征函数和处处稠密集补集的特征函数在 \(\mathcal{C}(\mathrm{X})\) 中生成一个布尔代数 A。随后他证明 A 的理想自然对应于空间 X 的闭子集，从而引入了一个后来特别富有成果的思想。

从1939年起，盖尔范德的观点为这些相对分散的结果打开了新视野。他的工作似乎部分受到维纳思想的推动，特别是其关于绝对收敛傅里叶级数的定理，以及彼得和外尔的方法。1939至1946年间，他与 M. 奈马克、D. 拉伊科夫和 G. 希洛夫合作，奠定了交换巴拿赫代数理论和星代数理论的基础（参见 {[}23{]}，第I卷，第169--400页）。特别地，他给出了维纳定理的经典证明——由 P. 莱维 {[}39{]} 推广——本书中可见该证明（I，第38页，例）。

在盖尔范德的方法中，与交换巴拿赫代数 A 相关的本质对象是 A 的极大理想空间 X。事实上，通过盖尔范德--马祖尔定理（I，第26页，推论2），A 的每个元素都在 X 上定义一个函数。众所周知，这一观点后来会在代数几何的发展中变得十分重要（参见 ÉHM，第146--148页）。这一方法本身受到了上文提到的斯通关于布尔代数的工作的启发。

似乎是 L. 卢米斯在1953年出版的课程 {[}41，第53页{]} 中首次强调 A 的特征标空间来阐述盖尔范德理论，并按本书所理解的意义定义盖尔范德变换。这一观点也许最清楚的优点是，特征标空间的拓扑性质直接来自巴拿赫空间对偶单位球在弱拓扑下的紧性。

从最早的工作起，盖尔范德就通过柯西积分在巴拿赫代数中构造了一元全纯函数演算，并由此推出：当极大理想空间不连通时，对合巴拿赫代数允许按环的乘积作非平凡分解。任意巴拿赫代数的情形直到1953年才得到处理，那时希洛夫在多元情形发展出一种函数演算。{[}68{]}。

至于星代数，盖尔范德和奈马克早在1942年就已证明它们可以实现为希尔伯特空间上的算子代数（参见 V，第442页定理2），它们很快成为大量深刻研究的对象；推动力来自群表示的应用、默里和冯·诺伊曼的工作以及量子力学的数学形式化。特别值得一提的是 I. 西格尔的贡献；他是系统研究星代数表示的先驱之一，这项研究与局部紧群的酉对偶研究相联系，我们将在几行后讨论后者。

尽管盖尔范德、奈马克、拉伊科夫和希洛夫在1946年以前已经奠定 C*-代数理论的基本基础，某些技术细节后来才得到改进。例如，R. 阿伦斯 {[}1{]} 说明了如何避免调用“希洛夫边界”（参见 I，第171页，习题26）的存在性，以证明盖尔范德变换是一个对合态射。此外，盖尔范德--奈马克关于含单位 C*-代数的理论起初把 \(1 + x^{*}x\) 对所有 \(x\) 可逆这一事实作为公理加入。然而他们猜想这个公理并非必要；这一断言的证明虽然初等却很精细，直到1952年前后才在 M. 福卡米亚 {[}21{]} 的工作、经 I. 卡普兰斯基 {[}62{]} 接续之后给出。

